\documentclass[10pt,oneside]{book}

\ifdefined\pdfinfoomitdate\pdfinfoomitdate=1\fi
\ifdefined\pdftrailerid\pdftrailerid{}\fi
\ifdefined\pdfsuppressptexinfo\pdfsuppressptexinfo=15\fi

\usepackage[T1]{fontenc}
\usepackage[utf8]{inputenc}
\usepackage[spanish,es-nodecimaldot]{babel}
\usepackage{amsmath,amssymb,mathrsfs}
\usepackage{aurical}
\usepackage[all]{xy}
\usepackage[a4paper,margin=22mm]{geometry}
\usepackage[colorlinks=true,linkcolor=blue,hypertexnames=false,unicode=true,pdfencoding=auto]{hyperref}
\hypersetup{
  pdftitle={Elementos de geometría algebraica III, segunda parte — edición española},
  pdfauthor={Alexander Grothendieck y Jean Dieudonné},
  pdfsubject={Estudio cohomológico de los haces coherentes: secciones 6 y 7}
}

\renewcommand{\thesection}{\arabic{section}}
\renewcommand{\thesubsection}{\thesection.\arabic{subsection}}
\newcommand{\sh}[1]{\mathscr{#1}}
% Sheaf Tor: script T followed by "or" in slanted Aurical, moved 2.2pt left.
\makeatletter
\newcommand*{\ega@sheafop}[3]{\sh{#1}#2\textup{\kern-2.2pt{\Fontauri\slshape #3}}}
\newcommand{\shTor}{\ega@sheafop{T}{}{or}}
\makeatother
\newcommand{\RR}{\mathrm R}
\newcommand{\HH}{\mathrm H}
\newcommand{\bb}[1]{\mathbf{#1}}
\newcommand{\cat}[1]{\mathtt{#1}}
\newcommand{\kres}{\kappa}
\newcommand{\vphi}{\phi}
\newcommand{\supertilde}{{\,\widetilde{\,}\,}}
\newcommand{\Spec}{\operatorname{Spec}}
\newcommand{\Proj}{\operatorname{Proj}}
\newcommand{\shProj}{\mathscr{P}\!roj}
\newcommand{\Hom}{\operatorname{Hom}}
\newcommand{\Tor}{\operatorname{Tor}}
\newcommand{\Ker}{\operatorname{Ker}}
\newcommand{\Coker}{\operatorname{Coker}}
\newcommand{\isoto}{\xrightarrow{\sim}}
\newcommand{\simto}{\xrightarrow{\sim}}
\newcommand{\shHom}{\mathscr{H}\!om}
\newcommand{\sref}[2][]{(#2)}
\newcommand{\oldpage}[2][]{\setcounter{footnote}{0}\marginpar{\scriptsize\textbf{#1}~|~#2}\ignorespaces}
\newenvironment{env}[1][]{\par\medskip\noindent\textbf{(#1)}\ }{\par}
\newenvironment{theorem}[1][]{\par\medskip\noindent\itshape Teorema (#1). ---\ }{\par}
\newenvironment{proposition}[1][]{\par\medskip\noindent\itshape Proposición (#1). ---\ }{\par}
\newenvironment{lemma}[1][]{\par\medskip\noindent\itshape Lema (#1). ---\ }{\par}
\newenvironment{corollary}[1][]{\par\medskip\noindent\itshape Corolario (#1). ---\ }{\par}
\newenvironment{remark}[1][]{\par\medskip\noindent\itshape Observación (#1). ---\ }{\par}
\newenvironment{remarks}[1][]{\par\medskip\noindent\itshape Observaciones (#1). ---\ }{\par}
\newenvironment{scholium}[1][]{\par\medskip\noindent\textit{Escolio (#1). ---\ }}{\par}
\newenvironment{definition}[1][]{\par\medskip\noindent\itshape Definición (#1). ---\ }{\par}
\newenvironment{proof}{\par\medskip\noindent\textit{Demostración. ---\ }}{\par}
\newenvironment{namedtheorem}[2][]{\par\medskip\noindent\itshape #2 (#1). ---\ }{\par}

\newcommand{\EGAUnit}[1]{\input{units/#1.es.tex}}

\setlength{\emergencystretch}{6em}
\begin{document}
% BEGIN RIGHTS AND AI TRANSLATION NOTICE
\begingroup
\small
\noindent\textbf{Derechos, disponibilidad y autoría}

Los derechos sobre las obras originales y sobre las contribuciones protegidas a estas traducciones y transcripciones derivadas corresponden a sus respectivos titulares. Todo aquello sobre lo que el editor tenga derechos en estas ediciones queda dedicado al dominio público bajo CC0 1.0 (https://creativecommons.org/publicdomain/zero/1.0/deed.es). Las menciones de derechos de terceros y las autorizaciones ya concedidas siguen siendo aplicables.

Estas ediciones se ofrecen gratuitamente, con fines no comerciales y exclusivamente educativos, bajo la presunción de que los respectivos titulares de derechos no se opondrían a que estas traducciones y transcripciones estuvieran disponibles gratuitamente junto a los originales franceses en NUMDAM (www.numdam.org/). Esta presunción no significa que se haya obtenido autorización; el acceso gratuito a los originales no autoriza por sí solo su traducción, adaptación o republicación.

Si usted es titular de derechos sobre algún material presentado aquí y se opone a su disponibilidad, puede abrir una incidencia en GitHub (github.com/KokunoYumeto/ega-fga-es/issues). El responsable de la publicación retirará el material afectado a petición del titular de los derechos.

La traducción al español está producida por IA. Las correcciones, la edición y la composición también están producidas por IA. El responsable de la publicación no tradujo el texto. Se conservan la autoría y las atribuciones de las obras originales y de las contribuciones heredadas; no se atribuyen a la IA. Los modelos documentados se identifican en cada edición.
\par
\endgroup
\clearpage
% END RIGHTS AND AI TRANSLATION NOTICE

\begin{titlepage}
  \centering
  {\Large A. Grothendieck y J. Dieudonné\par}
  \vfill
  {\Huge Elementos de geometría algebraica III\par}
  \vspace{1.5ex}
  {\LARGE Segunda parte\par}
  \vfill
  {\large Traducción española\par}
  \vspace{1ex}
  {\small Texto original: A. Grothendieck y J. Dieudonné, Publications Mathématiques de l’IHÉS, volumen 17 (1963)\par}
\end{titlepage}

\chapter*{CAPÍTULO III (continuación)\\[1.2ex]
\Large ESTUDIO COHOMOLÓGICO DE LOS HACES COHERENTES}

\setcounter{section}{5}
\setcounter{subsection}{0}
\setcounter{equation}{0}
\setcounter{footnote}{0}
\EGAUnit{III.6.1}
\EGAUnit{III.6.2}
\EGAUnit{III.6.3}
\EGAUnit{III.6.4}
\EGAUnit{III.6.5}
\EGAUnit{III.6.6}
\EGAUnit{III.6.7}
\EGAUnit{III.6.8}
\EGAUnit{III.6.9}
\EGAUnit{III.6.10}
\EGAUnit{III.7.1}
\EGAUnit{III.7.2}
\EGAUnit{III.7.3}
\EGAUnit{III.7.4}
\EGAUnit{III.7.5}
\EGAUnit{III.7.6}
\EGAUnit{III.7.7}
\EGAUnit{III.7.8}
\EGAUnit{III.7.9}
\EGAUnit{III.backmatter}
\end{document}
